
\documentclass[UTF8,zihao=-4]{ctexart}
\usepackage[a4paper,top=2.3cm,bottom=2.3cm,left=2.35cm,right=2.35cm]{geometry}
\usepackage{fontspec}
\usepackage{xeCJK}
\setmainfont{DejaVu Sans}
\setsansfont{DejaVu Sans}
\setmonofont{Noto Sans Mono}
\setCJKmainfont{Noto Sans CJK SC}
\setCJKsansfont{Noto Sans CJK SC}
\setCJKmonofont{Noto Sans CJK SC}
\usepackage{booktabs}
\usepackage{longtable}
\usepackage{array}
\usepackage{tabularx}
\usepackage{enumitem}
\usepackage{xcolor}
\usepackage{hyperref}
\usepackage{fancyhdr}
\usepackage{titlesec}
\usepackage{lastpage}
\usepackage{pifont}
\usepackage{amsmath}
\usepackage{amssymb}
\usepackage{makecell}
\hypersetup{colorlinks=true,linkcolor=black,urlcolor=blue,citecolor=black}
\definecolor{RuleBlue}{RGB}{25,74,135}
\definecolor{LightBlue}{RGB}{236,243,252}
\definecolor{LightGray}{RGB}{247,247,247}
\definecolor{DarkGray}{RGB}{80,80,80}
\definecolor{WarnOrange}{RGB}{172,103,0}
\definecolor{RejectRed}{RGB}{150,35,35}
\definecolor{PassGreen}{RGB}{38,118,82}
\titleformat{\section}{\Large\bfseries\color{RuleBlue}}{\thesection}{0.8em}{}
\titleformat{\subsection}{\large\bfseries}{\thesubsection}{0.8em}{}
\titleformat{\subsubsection}{\normalsize\bfseries}{\thesubsubsection}{0.8em}{}
\setlist[itemize]{topsep=3pt,itemsep=2pt,parsep=0pt,leftmargin=1.9em}
\setlist[enumerate]{topsep=3pt,itemsep=2pt,parsep=0pt,leftmargin=2.2em}
\newcolumntype{Y}{>{\raggedright\arraybackslash}X}
\newcolumntype{L}[1]{>{\raggedright\arraybackslash}p{#1}}
\newcommand{\Pass}{\textcolor{PassGreen}{\ding{51}}}
\newcommand{\Warn}{\textcolor{WarnOrange}{\ding{115}}}
\newcommand{\Reject}{\textcolor{RejectRed}{\ding{55}}}
\newcommand{\Code}[1]{\texttt{#1}}
\newcommand{\Term}[1]{\textbf{#1}}
\newenvironment{notebox}[1]{\vspace{0.5em}\noindent\colorbox{LightBlue}{\parbox{0.96\linewidth}{\textbf{#1}}}\par\vspace{-0.2em}\begin{quote}\small}{\end{quote}\vspace{0.3em}}
\newenvironment{decisionbox}[1]{\vspace{0.5em}\noindent\colorbox{LightGray}{\parbox{0.96\linewidth}{\textbf{#1}}}\par\vspace{-0.2em}\begin{quote}\small}{\end{quote}\vspace{0.3em}}

\hypersetup{pdftitle={Robot Demonstration Data Quality Framework System Design v1.0}, pdfauthor={ChatGPT for Lingji Qiwuzhi}}
\pagestyle{fancy}
\fancyhf{}
\lhead{RDDQF System Design v1.0}
\rhead{L3 v2 Roadmap}
\cfoot{第 \thepage\ 页}

\begin{document}

\begin{titlepage}
\centering
\vspace*{2.5cm}
{\Huge\bfseries 机器人 Demonstration 数据质量框架系统设计\par}
\vspace{0.35cm}
{\Large Robot Demonstration Data Quality Framework System Design\par}
\vspace{0.6cm}
{\LARGE v1.0\par}
\vspace{1.2cm}
\begin{tabular}{rl}
文档定位： & 覆盖 L3 v2 及后续文档体系的系统级设计说明 \\
适用系统： & 灵机启物机器人数采质检平台 \\
核心目标： & 从人工 QC 工具演进为机器人数据质量平台 \\
覆盖范围： & 标准、本体、指标、算法、实现、UI、报告、数据资产 \\
当前阶段： & 完成 Standard 与 Ontology 后的系统路线图 \\
\end{tabular}
\vfill
\begin{minipage}{0.86\linewidth}
\centering\large\bfseries
本系统设计的核心思想是：先建立训练数据质量标准，再建立质量本体，然后将指标、算法、代码、前端和报告全部挂载到同一套质量语义上。
\end{minipage}
\vfill
{\large 2026 年 6 月\par}
\end{titlepage}

\tableofcontents
\newpage

\section*{执行摘要}
\addcontentsline{toc}{section}{执行摘要}

本文档描述“机器人 Demonstration 数据质量框架”（Robot Demonstration Data Quality Framework, RDDQF）的系统级设计思路。它不是单个指标文档，也不是后端实现说明，而是覆盖后续所有设计文档的总纲。

系统目标是将当前 L3 从一个 episode 打开时临时计算的 Motion Score，升级为面向具身智能训练数据的数据质量平台能力。升级后的系统应具备：质量标准、质量本体、指标规范、算法规范、离线缓存、批次统计、前端质量画像、人工复核闭环、数据集发布报告和版本治理。

本文档提前定义后续需要编写的核心文档：Metric Specification、Algorithm Specification、Threshold Calibration Report、L3 Engine Implementation Specification、Manual QC UI Specification、Batch Quality Report Specification、Database Schema and Cache Design、Dataset Release Quality Protocol。每份文档的职责、输入、输出和与现有系统的关系均在本文档中说明。

\section{系统定位}

\subsection{当前系统现状}

当前平台已经具备完整的数据质检基本链路：MinIO 保存 raw 与 processed 文件，PostgreSQL 保存 episode 元数据和 QC 结果，FastAPI 后端提供 QC context 和 telemetry curve API，Vue 前端提供人工 QC 工作台。L3 v1 后端从 \Code{telemetry.npz} 读取 qpos、actions、qvel、effort、timestamps、sync validation 等字段，计算 8 个自动指标和一个 \Code{Q\_motion} 综合分。

该系统已经可以支持人工审核，但其质量语义仍偏向“运动执行质量”。下一阶段需要将它升级为“训练数据质量评价系统”。

\subsection{目标系统定位}

目标系统定义为：
\begin{quote}
面向机器人策略学习数据集构建的多层质量评价、质量追溯和数据资产筛选系统。
\end{quote}

它不是单纯的 QC 页面，也不是单纯的指标计算脚本，而是连接采集、审核、数据筛选、训练集构建和数据发布的质量基础设施。

\section{总体设计原则}

\subsection{标准先于指标}

系统不应从“能算哪些指标”出发，而应从“什么样的数据值得训练”出发。所有指标都必须能追溯到 Training Data Quality Standard 和 Quality Ontology。

\subsection{本体先于算法}

同一个质量节点可以有多个算法版本。例如 Smoothness 可以由 LDLJ、SPARC、EE jerk 或混合指标实现。本体稳定，算法可迭代。

\subsection{训练价值优先于执行精度}

L3 v2 的默认目标不是评价机器人是否精确执行遥操作命令，而是评价 demonstration 是否为模型提供高质量学习信号。Tracking Error 和 Effort 属于诊断类指标，而不是训练质量主指标。

\subsection{单条质量与数据集质量分离}

Episode 级指标回答“这条数据是否好”；Dataset 级指标回答“这一批数据是否覆盖充分、分布合理、冗余可控”。两者需要共享本体，但不能混成一个分数。

\section{四层文档体系}

\subsection{文档栈概览}

\begin{verbatim}
Layer 1: Training Data Quality Standard
         -> defines what high-quality training data means

Layer 2: Quality Ontology
         -> defines the quality tree and semantic categories

Layer 3: Metric Specification
         -> defines metrics under each ontology node

Layer 4: Algorithm and Implementation Specification
         -> defines formulas, APIs, code, cache, UI, reports
\end{verbatim}

\subsection{已完成文档}

\begin{longtable}{L{4cm} L{4cm} L{6cm}}
\toprule
文档 & 状态 & 作用 \\
\midrule
Training Data Quality Standard v1.0 & 已完成 & 定义训练数据质量的评价目标、哲学和准入原则。 \\
Robot Demonstration Data Quality Ontology v1.0 & 本次完成 & 定义 L3 v2 指标体系的质量本体树。 \\
\bottomrule
\end{longtable}

\subsection{后续核心文档}

\begin{longtable}{L{4.2cm} L{4.2cm} L{5.8cm}}
\toprule
文档 & 主要问题 & 产出 \\
\midrule
Metric Specification & 每个质量节点用什么指标表达 & 指标清单、定义、输入输出、解释逻辑。 \\
Algorithm Specification & 每个指标怎么算 & 公式、窗口、归一化、评分、timeline 规则。 \\
Threshold Calibration Report & Good/Warn/Bad 如何标定 & 分布统计、任务类型阈值、标注样本分析。 \\
L3 Engine Implementation Specification & 后端如何实现 & 模块结构、类设计、缓存、异步任务、API。 \\
Manual QC UI Specification & 前端如何展示质量画像 & 右侧面板、timeline、曲线联动、复核工作流。 \\
Batch Quality Report Specification & 批次级如何分析 & 分布、异常聚类、同类对比、数据集准入建议。 \\
Database Schema and Cache Design & 结果如何持久化 & quality dimension、metric、algorithm、result、版本表。 \\
Dataset Release Quality Protocol & 如何发布训练数据集 & 数据集切分、质量报告、版本号、可追溯记录。 \\
\bottomrule
\end{longtable}

\section{目标系统架构}

\subsection{逻辑架构}

\begin{verbatim}
MinIO processed data
      |
      v
Telemetry Loader / Metadata Loader
      |
      v
Feature Extraction Layer
      |
      v
Ontology-aware Metric Engine
      |
      +--> Episode Quality Profile
      +--> Timeline Segments
      +--> Diagnostic Signals
      |
      v
Quality Result Store / Cache
      |
      +--> Manual QC UI
      +--> Batch Quality Report
      +--> Dataset Release Report
      +--> Threshold Calibration Pipeline
\end{verbatim}

\subsection{核心模块}

\begin{longtable}{L{4cm} L{5cm} L{5cm}}
\toprule
模块 & 职责 & 备注 \\
\midrule
Data Access Layer & 从 MinIO/PostgreSQL 读取 telemetry、manifest、metadata、视频索引 & 避免多个 API 重复读取同一对象。 \\
Feature Extraction Layer & 生成 arm/hand/EE/time/sync 等基础特征 & 与具体指标解耦，可缓存。 \\
Quality Ontology Registry & 保存质量维度、指标归属和版本关系 & 支持 UI 分组与报告生成。 \\
Metric Engine & 计算 episode 级指标 & 按本体维度组织输出。 \\
Timeline Generator & 生成异常、低信息、关键动作、复核片段 & 不再只输出异常段。 \\
Quality Result Store & 持久化指标结果和算法版本 & 支持离线缓存和历史对比。 \\
Calibration Pipeline & 基于 approved samples 做统计标定 & 从经验阈值升级为数据驱动阈值。 \\
Report Generator & 生成 episode/batch/dataset 级报告 & 服务数据资产发布。 \\
Manual QC UI & 人工复核、解释、确认和反馈 & 提供闭环标注。 \\
\bottomrule
\end{longtable}

\section{从 L3 v1 到 L3 v2 的迁移路线}

\subsection{保留但重解释的内容}

L3 v1 中并非所有工作都需要推翻。已有的数据读取、关节维度识别、timeline 合并、前端曲线展示、settings 配置页等仍有价值。需要改变的是指标语义和评分结构。

\begin{longtable}{L{3.5cm} L{4.6cm} L{5.5cm}}
\toprule
L3 v1 元素 & L3 v2 处理 & 原因 \\
\midrule
LDLJ & 保留并归入 Smoothness & 可表达示范平滑性。 \\
Dead/Static & 合并为低信息/低活动片段证据 & 避免重复惩罚。 \\
Saturation & 语义化处理 & 手部开合极限可能正常。 \\
Jitter/Sync & 保留为 Data Integrity 硬门控 & 直接影响 observation-action 对齐。 \\
Tracking Error & 降级为 Execution Diagnostics & 评价控制器，不直接评价训练价值。 \\
Effort & 降级为 Execution Diagnostics & 辅助排查负载和碰撞风险。 \\
Chatter & 保留并升级为 Finger Command Quality & 对灵巧手训练价值重要。 \\
Q\_motion & 替换为 Quality Profile & 单一运动分不足以表达训练质量。 \\
\bottomrule
\end{longtable}

\subsection{新增能力}

L3 v2 应新增以下能力：
\begin{itemize}
  \item EE pose 相关质量：末端轨迹平滑度、末端路径效率、末端稳定性。
  \item 任务阶段感知：接近、抓取、移动、放置、完成等阶段片段。
  \item 低信息片段识别：将无效等待、空转和真实关键停顿区分开。
  \item 训练标签质量：action label 的平滑性、动态范围、跳变、手部协同。
  \item 质量画像输出：不再只输出一个分数，而是输出多维 quality profile。
  \item 离线缓存：避免每次打开 manual QC 都重新计算。
  \item 批次对比：同类任务分布、离群样本、冗余样本、质量分布。
\end{itemize}

\section{Metric Specification 文档预案}

\subsection{文档目标}

Metric Specification 是下一份应编写的核心文档。它回答：对于 Ontology 中的每个质量节点，需要哪些指标表达？每个指标为什么存在？依赖哪些数据？输出什么？怎么解释？

\subsection{建议结构}

\begin{enumerate}
  \item 指标命名规范。
  \item 指标分层：core、contextual、diagnostic、dataset-level。
  \item 每个质量维度下的候选指标表。
  \item 现有指标迁移清单。
  \item 新增指标优先级。
  \item 指标输出 schema。
  \item Timeline 支持规则。
  \item UI 展示建议。
\end{enumerate}

\subsection{候选指标一级清单}

\begin{longtable}{L{4cm} L{5.5cm} L{4.8cm}}
\toprule
质量维度 & 候选指标方向 & 优先级 \\
\midrule
Learnability & low information ratio, action-state activity, effective motion ratio & P0 \\
Demonstration Behavior & joint smoothness, EE smoothness, trajectory efficiency, correction burden & P0/P1 \\
Action Label Quality & action jump rate, command smoothness, hand chatter, saturation semantics & P0/P1 \\
State Trajectory & qpos continuity, EE continuity, workspace path length, pose stability & P1 \\
Data Integrity & timestamp jitter, sync invalid ratio, frame consistency, metadata consistency & P0 \\
Execution Diagnostics & command tracking, effort p95, IMU shock & P1/P2 \\
Dataset-Level & task distribution, embedding diversity, duplicate ratio, outlier score & P2 \\
\bottomrule
\end{longtable}

\section{Algorithm Specification 文档预案}

\subsection{文档目标}

Algorithm Specification 在 Metric Specification 之后编写。它才开始给出公式、窗口、阈值映射和 timeline 生成规则。

\subsection{每个算法条目应包含}

\begin{itemize}
  \item Metric ID 与所属 ontology node。
  \item 输入字段和预处理逻辑。
  \item 计算公式或伪代码。
  \item 输出 value、score、severity、confidence。
  \item 时间复杂度和内存复杂度。
  \item Timeline segment 生成规则。
  \item 已知误报场景和任务类型适配策略。
  \item 标定所需样本和验证方法。
\end{itemize}

\section{Threshold Calibration 文档预案}

\subsection{为什么需要单独文档}

当前 L3 v1 的阈值主要是经验值。L3 v2 必须从 approved samples 和 failed samples 的实际分布中标定阈值。阈值不是算法的一部分，而是数据集和任务类型相关的配置资产。

\subsection{标定流程}

\begin{enumerate}
  \item 收集每个任务类型至少 200 条已人工确认样本。
  \item 计算候选指标分布。
  \item 对 approved、warning、rejected 数据进行分布对比。
  \item 选取 percentile 或 robust statistics 作为初始阈值。
  \item 人工复核边界样本。
  \item 建立任务类型阈值 profile。
  \item 持续记录阈值版本与效果。
\end{enumerate}

\section{L3 Engine Implementation 文档预案}

\subsection{后端模块建议}

\begin{verbatim}
backend/app/services/l3_v2/
|-- loaders.py              # telemetry, manifest, metadata loading
|-- features.py             # reusable feature extraction
|-- ontology.py             # quality dimension registry
|-- metrics/
|   |-- learnability.py
|   |-- behavior.py
|   |-- action_label.py
|   |-- trajectory.py
|   |-- integrity.py
|   `-- diagnostics.py
|-- scoring.py              # quality profile aggregation
|-- timeline.py             # segment generation
|-- cache.py                # result cache
`-- engine.py               # orchestration
\end{verbatim}

\subsection{API 输出建议}

\begin{verbatim}
{
  "episodeId": "...",
  "qualityProfile": {
    "learnability": {...},
    "demonstrationBehavior": {...},
    "actionLabelQuality": {...},
    "dataIntegrity": {...},
    "executionDiagnostics": {...}
  },
  "metricCards": [...],
  "timelineSegments": [...],
  "diagnostics": [...],
  "algorithmVersions": {...},
  "confidence": {...}
}
\end{verbatim}

\section{Manual QC UI 文档预案}

\subsection{UI 目标}

前端应从“展示一堆指标卡片”升级为“解释这条数据是否值得训练”。右侧边栏应展示 quality profile，曲线和 timeline 应围绕质量维度联动。

\subsection{建议 UI 结构}

\begin{itemize}
  \item 顶部：episode 基本信息、任务类型、当前训练质量总览。
  \item 主区：多视角视频、轨迹曲线、异常/关键片段 timeline。
  \item 右侧：Training Quality Profile，不同维度独立展示。
  \item 诊断区：Tracking、Effort、IMU 等辅助信息折叠展示。
  \item 人工复核：Pass/Fail 之外增加“适合训练/仅供诊断/需人工裁剪/不进入训练集”等数据资产状态。
\end{itemize}

\section{Batch Quality Report 文档预案}

\subsection{报告目标}

批次报告不只是列出每条数据分数，而是回答：这一批数据能不能构成可用训练集？哪些任务覆盖不足？哪些 episode 是离群样本？哪些样本重复？哪些指标分布异常？

\subsection{核心内容}

\begin{itemize}
  \item 任务类型分布与采集规模。
  \item 各质量维度分布图。
  \item 按任务类型的质量分位数。
  \item 低质量原因聚类。
  \item 高价值样本推荐。
  \item 重复/冗余 episode 检测。
  \item 数据集发布建议。
\end{itemize}

\section{Database Schema and Cache Design 文档预案}

\subsection{核心表设计方向}

\begin{longtable}{L{4cm} L{8.5cm}}
\toprule
表名 & 职责 \\
\midrule
\Code{quality\_dimensions} & 存储 ontology 节点。 \\
\Code{metric\_definitions} & 存储指标语义、所属维度、核心/诊断/批次类型。 \\
\Code{metric\_algorithms} & 存储算法版本、参数 schema 和实现版本。 \\
\Code{episode\_metric\_results} & 存储 episode 级指标结果。 \\
\Code{episode\_quality\_profiles} & 存储 episode 级多维质量画像。 \\
\Code{quality\_segments} & 存储 timeline segments。 \\
\Code{threshold\_profiles} & 存储任务类型阈值配置。 \\
\Code{batch\_quality\_reports} & 存储批次级分析结果。 \\
\bottomrule
\end{longtable}

\subsection{缓存原则}

同一 telemetry 文件、同一算法版本、同一阈值 profile 下的结果应可复用。缓存 key 至少应包含 episode id、processed object version、algorithm version、config version。

\section{Dataset Release Quality Protocol 文档预案}

\subsection{数据资产状态}

建议为每条 episode 增加数据资产状态：
\begin{itemize}
  \item \Code{train\_ready}: 可直接进入训练集。
  \item \Code{train\_with\_caution}: 可训练，但存在低风险提示。
  \item \Code{needs\_trim}: 需要裁剪低质量片段后训练。
  \item \Code{diagnostic\_only}: 不适合作为正样本，但可用于诊断或失败分析。
  \item \Code{reject}: 不进入数据资产库。
\end{itemize}

\subsection{发布报告内容}

每个训练数据集版本应附带质量报告，包括样本数量、任务分布、质量维度分布、剔除原因、阈值版本、指标算法版本、人工 QC 版本和可追溯链接。

\section{实施路线图}

\subsection{Phase 0：设计基线}

\begin{itemize}
  \item 已完成 Training Data Quality Standard。
  \item 已完成 Quality Ontology。
  \item 完成本系统设计文档。
\end{itemize}

\subsection{Phase 1：Metric Specification}

产出完整指标体系，不写代码。目标是确定 P0/P1/P2 指标清单，以及每个指标的本体挂载位置。

\subsection{Phase 2：Algorithm Specification}

对 P0 指标给出公式、窗口、score 映射、timeline 规则和误报场景。

\subsection{Phase 3：后端 L3 v2 原型}

在不破坏 L3 v1 的前提下实现 L3 v2 engine，输出 quality profile。优先使用离线缓存。

\subsection{Phase 4：前端质量画像改造}

将 manual QC 页面从指标卡片列表升级为 quality profile，支持点击 timeline 跳转视频，支持曲线异常高亮。

\subsection{Phase 5：批次统计与阈值标定}

基于 200+ approved samples 开始任务类型阈值标定，并生成 batch quality report。

\subsection{Phase 6：数据资产发布协议}

将 QC 结果用于训练集筛选、数据集版本发布和训练前数据报告。

\section{近期下一步建议}

基于当前进度，下一步应立即进入 \Term{Metric Specification}。建议优先完成以下输出：

\begin{enumerate}
  \item P0 指标候选表：必须支持训练数据准入。
  \item P1 指标候选表：提升质量解释能力。
  \item Diagnostic 指标表：Tracking、Effort、IMU 等辅助诊断。
  \item 现有 L3 v1 指标迁移表。
  \item 新旧 API 输出对比。
  \item 前端质量画像初稿。
\end{enumerate}

\section{结论}

RDDQF 的核心不是增加更多指标，而是建立一条从质量标准到本体、从本体到指标、从指标到算法、从算法到平台能力的稳定链路。只有这样，L3 v2 才不会再次退化为经验指标堆叠，而能成为机器人数据资产平台的质量基础设施。

\end{document}
